% Generated draft: requires mathematical and source review.
\documentclass[11pt]{article}
\usepackage[margin=25mm]{geometry}
\usepackage{fontspec}
\setmainfont{Latin Modern Roman}
\usepackage{amsmath,amssymb,mathtools,unicode-math}
\setmathfont{Latin Modern Math}
\usepackage{enumitem,needspace,xurl,hyperref}
\setlength{\parindent}{0pt}
\setlength{\parskip}{5pt}
\setlength{\emergencystretch}{3em}
\newcommand{\N}{\mathbb{N}}
\newcommand{\Z}{\mathbb{Z}}
\newcommand{\Q}{\mathbb{Q}}
\newcommand{\R}{\mathbb{R}}
\newcommand{\C}{\mathbb{C}}
\newcommand{\F}{\mathbb{F}}
\newcommand{\A}{\mathbb{A}}
\newcommand{\PP}{\mathbb{P}}
\newcommand{\E}{\mathbb{E}}
\newcommand{\eps}{\varepsilon}
\newcommand{\dd}{\,\mathrm{d}}
\begin{document}
\section*{UnsolvedMath ID: 30002953}
\textbf{Source problem code:} catalogue entry
\subsection{Definitions and mathematical statement}
The supplied statement uses cumulative distribution functions and sums without explicitly specifying the underlying state space. The following formulation uses the conventional interpretation of real-valued random variables and a finite positive integer \(n\); this interpretation requires human review against the cited report.

Let \(n\geq 1\), and let \(F_1,\ldots,F_n\) be cumulative distribution functions of Borel probability measures \(\mu_1,\ldots,\mu_n\) on \(\R\), so that \(F_i(t)=\mu_i(({-\infty},t])\). Define the set of all couplings of these marginals by
\[
\Pi(\mu_1,\ldots,\mu_n)
=
\left\{\pi:\ \pi\text{ is a Borel probability measure on }\R^n,\quad
\pi\bigl(\{x\in\R^n:x_i\in A\}\bigr)=\mu_i(A)
\text{ for every Borel }A\subseteq\R\text{ and every }i\right\}.
\]
For \(\pi\in\Pi(\mu_1,\ldots,\mu_n)\), let
\[
H_\pi(t)=\pi\left(\left\{x\in\R^n:\sum_{i=1}^n x_i\leq t\right\}\right),
\qquad t\in\R,
\]
and define the attainable-sums set
\[
D_n(F_1,\ldots,F_n)
=\{H_\pi:\pi\in\Pi(\mu_1,\ldots,\mu_n)\}.
\]
Thus all dependence structures are allowed. No moment, continuity, bounded-support, or independence assumptions are supplied.

\textbf{Requested characterization.} For arbitrary prescribed marginals \(F_1,\ldots,F_n\), give a direct characterization of \(D_n(F_1,\ldots,F_n)\), equivalently necessary and sufficient conditions on a candidate cumulative distribution function \(H\) for
\[
H\in D_n(F_1,\ldots,F_n).
\]
In addition, find a general analytical method to decide whether these marginals are jointly mixable, meaning that
\[
\exists c\in\R\ \exists\pi\in\Pi(\mu_1,\ldots,\mu_n):\quad
\pi\left(\left\{x\in\R^n:\sum_{i=1}^n x_i=c\right\}\right)=1.
\]
Equivalently, writing \(H_c(t)=0\) for \(t<c\) and \(H_c(t)=1\) for \(t\geq c\), joint mixability means
\[
\exists c\in\R:\quad H_c\in D_n(F_1,\ldots,F_n).
\]
Here the source's phrase ``contains a point mass'' means that the set contains the cumulative distribution function of a point-mass probability measure. The constant \(c\) is not prescribed in the supplied problem. The source does not specify what form of characterization or analytical decision method is required, or how arbitrary marginal distributions are to be represented as inputs; no computational decidability requirement is inferred.
\subsection{Short English statement}
Given marginal distributions, describe every distribution that their sum can have when dependence is unrestricted. Also seek a general analytical criterion for whether a coupling makes the sum equal to some constant almost surely. The real-valued interpretation and the intended meaning of an analytical decision method require review. The supplied record labels the problem open; that status is dated curation, not an independently verified conclusion.
\subsection{Sources}
\begin{itemize}
\item Local registry: \texttt{data/unsolvedmath/problems.json}, supplied complete source record for permanent UnsolvedMath ID \(30002953\), entitled ``Attainable Sum Distributions and Joint Mixability''; exact source problem code \texttt{catalogue entry}. The record belongs to the Oberwolfach Reports Open Problems collection, records proposed year 2015, and supplies no named proposer.
\item Supplied source citation: \emph{The Mathematics and Statistics of Quantitative Risk Management} (2016), Oberwolfach Reports, DOI 10.4171/owr/2015/42, pp. 2481--2532. Supplied source link: \url{https://doi.org/10.4171/owr/2015/42}. The background describes Oberwolfach Reports as documenting mathematical workshops and their open-problem discussions.
\item Supplied literature reference: ``Oberwolfach Report 42, 2015,'' \url{https://doi.org/10.4171/owr/2015/42}. The exact supplied literature-check date is \texttt{2026-08-21}. This is supplied dated curation, not independently verified here. Its assessment states that the report explicitly identifies the absence of a direct characterization and a general analytical verification method for \(D_n\), and that no complete general characterization was identified. The supplied background requests independent verification of time-sensitive status and cited claims.
\item The registry supplies status \texttt{open} and statement status \texttt{corrected\_verified}. Its verification note reports a manual semantic, typographic, and LaTeX review of the cleanup. These are supplied record annotations; this statement-only draft does not independently establish mathematical correctness or current literature status.
\end{itemize}
\end{document}
